% Fijiisland
% 2021-12-18
\chapter{序言}\label{PREFACE}

Rust是一门为系统编程而生的语言。

\medskip
然而在如今，绝大多数正在工作的程序员对于系统编程并不熟悉，那么便需要对其进行一些解释了，因为至今它仍是我们做任何事的基础。

\medskip
你合上了笔记本电脑。操作系统会感知到，并中止所有正在运行的程序，熄灭屏幕，最后让电脑进入睡眠。过了一会儿，你又打开了
电脑，屏幕与其它组件再次被启动，并且每个程序均能从它中断的地方恢复过来，我们认为这是理所当然的事。但这都是系统程序员
写了大量代码而得到的成果。

\medskip
系统编程一般被用到这些地方：
\begin{itemize}
    \item 操作系统
    \item 各种设备驱动
    \item 文件系统
    \item 数据库
    \item 运行在廉价设备、亦或是极度可靠的设备上的代码
    \item 密码学
    \item 媒体解码器（负责读取及写入音频、视频或者图像文件的软件）
    \item 媒体处理（例如，语音识别或照片编辑软件）
    \item 内存管理（例如，实现一个垃圾回收器）
    \item 文本渲染（将文本和字体转变为像素）
    \item 实现高级编程语言（比如 JavaScript 和 Python）
    \item 网络
    \item 虚拟化技术以及软件容器
    \item 科学模拟
    \item 游戏
\end{itemize}
简单来说，系统编程是一种资源受限的程序设计，它面向每个字节(byte)和每个CPU周期(cycle)来编程。

\medskip
涉及支持一个基础应用程序运行的系统代码的数量会相当惊人。

%%%%%%%%%%%%%%%%%%%%%%%%%%%%%%%%%%%%%%%%%%%%%%%%%%%%%%%%%%%%%%%%%%%%%%%%%%%%%%%%%%%%%%%%

\mksec{适合阅读本书的人群}
如果你已是一位系统程序员并且准备好迎接C++的替代品，那么本书适合你。若你是对于任何编程语言均有经验的开发者，无论是
C\#，Java，Python，JavaScript或者是其它东西，那么本书也适合你。

\medskip
但是，你不能光学习Rust。要想最大化地利用该门语言，你还需获得一些系统编程的经验。我们建议你在阅读本书的同时用Rust
实现一些有关系统编程的项目。试着去做一些你从来没有构建过的，能利用好Rust的迅捷、并发以及安全性的东西。本序言开端列
出的那些议题应该能为你带来一些灵感。

%%%%%%%%%%%%%%%%%%%%%%%%%%%%%%%%%%%%%%%%%%%%%%%%%%%%%%%%%%%%%%%%%%%%%%%%%%%%%%%%%%%%%%%%

\mksec{我们编写本书的原因}
我们想编写一本我们自己当初学习Rust时便希望拥有的书。我们的目标是开门见山地，将Rust重要且新的概念整理出来，进而清晰
而深入地表述它们，以便减少“通过试错来学习”这样的现象\textbf{（待修订）}。

%%%%%%%%%%%%%%%%%%%%%%%%%%%%%%%%%%%%%%%%%%%%%%%%%%%%%%%%%%%%%%%%%%%%%%%%%%%%%%%%%%%%%%%%

\mksec{本书导航}
本书的前两章介绍了Rust，并在我们移步至\mkref{CHP3}{第三章}的基础数据类型前提供了一段简短的总览。有关所有权和引用的概念位于\mkref{CHP4}{第四章}与
\mkref{CHP5}{第五章}。我们建议按顺序通读本书前五章。

\medskip
第\mkref{CHP6}{六}至第\mkref{CHP10}{十}章涵盖了该门语言的基础知识：表达式（\mkref{CHP6}{第六章}）、错误处理（\mkref{CHP7}{第七章}）、包和模
块（\mkref{CHP8}{第八章}）、结构体（\mkref{CHP9}{第九章}）以及枚举和模式（\mkref{CHP10}{第十章}）。对这些章进行跳跃式阅读没问题，但一定别跳过错
误处理这章，相信我们。

\medskip
\mkref{CHP11}{第十一章}涵盖了特型和泛型，也是最后两个你需要了解的重要概念。特型类似于Java或C\#里的接口。同时它们作为一种主要方式，使得Rust能够允许你
将你的自定义类型整合进该语言本身。\mkref{CHP12}{第十二章}展示了特型是如何支持运算符重载的，另外\mkref{CHP13}{第十三章}涵盖了更多实用特型。

\medskip
了解特型和泛型之后便可解锁本书的余下部分。第\mkref{CHP14}{十四}、\mkref{CHP15}{十五}章涵盖了闭包和迭代器这两个你绝对不想错过的强大工具。对于剩余
章节你可以按自由顺序进行阅读，或者根据需求对它们进行深度挖掘。这些章节涵盖了该门语言的剩余部分：集合（\mkref{CHP16}{第十六章}）、字符串和文本
（\mkref{CHP17}{第十七章}）、输入输出（\mkref{CHP18}{第十八章}）、并发（\mkref{CHP19}{第十九章}）、异步编码（\mkref{CHP20}{第二十章}）、
宏（\mkref{CHP21}{第二十一章}）、unsafe编码（\mkref{CHP22}{第二十二章}）以及如何在其它语言中调用Rust函数（\mkref{CHP23}{第二十三章}）。

%%%%%%%%%%%%%%%%%%%%%%%%%%%%%%%%%%%%%%%%%%%%%%%%%%%%%%%%%%%%%%%%%%%%%%%%%%%%%%%%%%%%%%%%

\mksec{本书中的惯例用法}
以下排版惯例被用于本书中：

\medskip
意大利斜体 --- \textit{Italic}
    \begin{quote}
        表示新术语、网址、电子邮件地址、文件名以及文件拓展名。
    \end{quote}
\medskip
等宽体 --- \texttt{Constant width}
    \begin{quote}
        用来列出程序代码，也用作在段内引用如变量、函数名、数据库、数据类型、环境变量、声明式以及关键字这样的程序元素。
    \end{quote}
\medskip
等宽粗体 --- \textbftt{Constant width bold}
    \begin{quote}
        表示需由用户键入的命令或其它文本。
    \end{quote}
\medskip
等宽斜体 --- \textittt{Constant width italic}
    \begin{quote}
        表示需要被替换的文本，由来自用户的值或由上下文确定的值提供。
    \end{quote}
\mknote{该框体代表一个一般批注。}

%%%%%%%%%%%%%%%%%%%%%%%%%%%%%%%%%%%%%%%%%%%%%%%%%%%%%%%%%%%%%%%%%%%%%%%%%%%%%%%%%%%%%%%%

\mksec{使用示例代码}
补充材料（示例代码、编程练习等）可于下方网址下载

\smallskip
\reditxt{https://github.com/ProgrammingRust.}

\smallskip
本书是为了帮你完成工作的。一般来说，若示例代码由本书提供，你可以将它用于你的项目和文档中。若非大量地复制本书中的代码，那么
你不需要联系我们以取得授权。比如说，写一个会用到本书中几个代码块的程序，不需要取得授权；售卖或分发O'Reilly书籍中的实例则
需要取得授权；回复问题时提及本书并引用示例代码不需要取得授权；将本书中大量代码纳入你的产品文档中也不需要取得授权。

\medskip
我们不强制要求，但同时欢迎你在进行引用时注明出处。出处通常包含书名、作者、出版商以及书籍ISBN号。例如：
\textit{"Programming Rust,Second Edition by Jim Blandy, Jason Orendorff, and Leonora F.S. Tindall(O'Reilly). Copyright 2%
021 JimBlandy, Leonora F.S. Tindall, and Jason Orendorff, 978-1-492-05259-3."}

\medskip
若你认为你对于示例代码的使用脱离了合理范畴，或是不属于上述的授权论列，请尽情通过邮箱联系我们：

\smallskip
\reditxt{permissions@oreilly.com}
\smallskip

%%%%%%%%%%%%%%%%%%%%%%%%%%%%%%%%%%%%%%%%%%%%%%%%%%%%%%%%%%%%%%%%%%%%%%%%%%%%%%%%%%%%%%%%

\mksec{O'Reilly在线学习}
\mknote{进40年来，O'Reilly Media提供技术与商业培训、知识、以及卓越见解来帮助众多公司取得成功。}

\medskip
我们拥有独一无二的专家和革新者组成的庞大网络，他们通过图书、文章、会议和我们的在线学习平台分享他们的知识和经验。O'Reilly的
在线学习平台允许你按需访问现场培训课程、深入的学习路径、交互式编程环境，以及O'Reilly和200多家其他出版商提供的大量文本和视
频资源。对于更多信息，请访问\reditxt{http://oreilly.com}。

%%%%%%%%%%%%%%%%%%%%%%%%%%%%%%%%%%%%%%%%%%%%%%%%%%%%%%%%%%%%%%%%%%%%%%%%%%%%%%%%%%%%%%%%

\mksec{联系我们}
与本书有关的评论和问题，请发给出版社：
    \begin{quote}
        O'Reilly Media, Inc.

        \smallskip
        1005 Gravenstein Highway North

        \smallskip
        Sebastopol, CA 95472

        \smallskip
        800-998-9938（美国或加拿大境内）

        \smallskip
        707-829-0515（国际或本地）

        \smallskip
        707-829-0104（传真）
    \end{quote}
欢迎访问本书主页，我们列出了勘误表、示例以及其它额外信息。请访问\reditxt{https://oreil.ly/programming-rust-2e}。

\medskip
通过发送电邮至\reditxt{bookquestions@oreilly.com}以提出关于本书的评论或者技术问题。

\medskip
访问\reditxt{http://www.oreilly.com}获得更多有关我们的书籍和课程的信息。

\medskip
在Facebook上找到我们: \reditxt{http://facebook.com/oreilly}

\medskip
在Twitter上关注我们: \reditxt{http://twitter.com/oreillymedia}

\medskip
在Youtube上关注我们: \reditxt{http://youtube.com/oreillymedia}

%%%%%%%%%%%%%%%%%%%%%%%%%%%%%%%%%%%%%%%%%%%%%%%%%%%%%%%%%%%%%%%%%%%%%%%%%%%%%%%%%%%%%%%%

\mksec{致谢}
你手里所持的本书很大程度上得益于我们官方技术审稿人的关注: Brain Anderson, Carol Nichols, 和Erik Nordin; 以及我们的
译者: Hidemoto Nakada（中田 秀基）（日本语），以及Adam Bochenek和Krzysztof Sawka（波兰语）。

\medskip
许多其它非官方审稿人阅读了本书初稿并提供了宝贵的反馈。我们想对于Eddy Bruel, Nick Fitzgerald, Graydon Hoare, Michael
 Kelly, Jeffrey Lim, Jakob Olesen, Gian-Carlo Pascutto, Larry Rabinowitz, Jaroslav Šnajdr, Joe Walker, 
以及Yoshua Wuyts所提出的周到意见表示感谢。我们特别要感谢Jeff Walden和Nicolas Pierron，他们抽出时间对几乎整本书进行了审
校。就像任何企业编程项目，一本编程书籍伴随着高质量的疏漏反馈而茁壮成长，谢谢你们。

\medskip
Mozilla对于Jim和Jason在此项目上的工作极为包容，即使这并非我们工作职责所在，而且还会占用一部分工作时间。非常感谢Dave Camp, 
Naveed Ihsanullah, Tom Tromey和Joe Walker这几位领导的支持。他们都在为Mozilla的长远着想，我们希望这份成果能够达到他们
的期望。

\medskip
我们还想感谢所有帮助实现该项目的O'Reilly工作人员，特别感谢我们极耐心的Jeff Bleiel和Brian MacDonald编辑，以及我们的策划
编辑Zan McQuade。

\medskip
最重要的是，衷心感谢我们的家人，感谢他们坚定不移的爱、热情和耐心。